\documentclass[12pt]{amsart}
\pdfoutput=1
\usepackage{latexsym, amsmath, amssymb, amsthm, bm}
\usepackage{calrsfs}
\usepackage{subfig}

\usepackage[mathscr]{eucal}
\usepackage{mathtools}
\usepackage{paralist}
\usepackage[alphabetic]{amsrefs}
\usepackage[T1]{fontenc}
\usepackage{mathptmx}
\usepackage{microtype}
\usepackage{tikz}
\renewcommand{\baselinestretch}{1.05}
\usepackage[colorlinks=true,linkcolor=blue,urlcolor=blue, citecolor=blue]%
    {hyperref}
\linespread{1.06}
%%--LAYOUT--------------------------------------------------------------------
\usepackage[paperwidth=8.5in, paperheight=11in, inner=2.5cm, outer=2.5cm,% 
    top=2.5cm, bottom=2.5cm]{geometry}
%%--OTHER ENVIRONMENTS--------------------------------------------------------
\newtheorem{lemma}{Lemma}[section]
\newtheorem{theorem}[lemma]{Teorema}
\newtheorem{corollary}[lemma]{Corollary}
\newtheorem{proposition}[lemma]{Proposition}
\theoremstyle{definition}
\newtheorem{definition}[lemma]{Definition}
\newtheorem{remark}[lemma]{Remark}
\newtheorem{examplex}[lemma]{Example}
\newtheorem{procedure}[lemma]{Procedure}

\newenvironment{exm}
  {\pushQED{\qed}\renewcommand{\qedsymbol}{$\diamond$}\examplex}
  {\popQED\endexamplex}
\newtheorem{question}[lemma]{Question}
\newtheorem{conjecture}[lemma]{Conjecture}
\newtheorem{Construction}[lemma]{Construction}
\newtheorem{algorithm}[lemma]{Algorithm}
%\theoremstyle{remark}
\newcommand{\define}[1]{{\bfseries\itshape #1}}
%%--MATH----------------------------------------------------------------------
\newcommand{\relphantom}[1]{\mathrel{\phantom{#1}}}
\newcommand{\abs}[1]{\ensuremath{\left| #1 \right|}}
%\newcommand{\norm}[1]{\ensuremath{\left\| #1 \right\|}}
\newcommand{\ideal}[1]{\ensuremath{\left\langle #1 \right\rangle}}
\renewcommand{\geq}{\geqslant}
\renewcommand{\leq}{\leqslant}
\newcommand{\coloneq}{\ensuremath{\mathrel{\mathop :}=}}
\newcommand{\transpose}{\ensuremath{\textsf{\upshape T}}} 
%% Blackboard bolds symbols
\renewcommand{\AA}{\ensuremath{\mathbb{A}}} 
\newcommand{\CC}{\ensuremath{\mathbb{C}}} 
\newcommand{\kk}{\ensuremath{\Bbbk}} 
\newcommand{\NN}{\ensuremath{\mathbb{N}}}
\newcommand{\PP}{\ensuremath{\mathbb{P}}} 
\newcommand{\QQ}{\ensuremath{\mathbb{Q}}} 
\newcommand{\RR}{\ensuremath{\mathbb{R}}} 
\newcommand{\EE}{\ensuremath{\mathbb{E}}} 

\renewcommand{\P}{\ensuremath{P}} 
\newcommand{\Q}{\ensuremath{Q}} 

%\renewcommand{\SS}{\ensuremath{\mathbb{S}}} 
\newcommand{\ZZ}{\ensuremath{\mathbb{Z}}} 
%% Script symbols
\newcommand{\sC}{\ensuremath{\mathcal{C}}}
\newcommand{\sE}{\ensuremath{\mathcal{E}}}
\newcommand{\sF}{\ensuremath{\mathcal{F}}}
\newcommand{\calF}{\ensuremath{\mathcal{F}}}
\newcommand{\calL}{\ensuremath{\mathcal{L}}}

\newcommand{\sH}{\ensuremath{\mathcal{H}}}
\newcommand{\sL}{\ensuremath{\mathcal{L}}}
\newcommand{\sO}{\ensuremath{\mathcal{O}}}
\newcommand{\diam}{\operatorname{diam}}
\newcommand{\Osh}{{\mathcal O}}
\newcommand{\Ish}{{\mathcal I}}
\newcommand{\calI}{{\mathcal I}}

%% Operators
\DeclareMathOperator{\codim}{codim}
\DeclareMathOperator{\coker}{Coker}
\DeclareMathOperator{\conv}{conv}
\DeclareMathOperator{\Cox}{Cox}
\DeclareMathOperator{\depth}{depth}
\DeclareMathOperator{\gon}{gon}
\DeclareMathOperator{\Grass}{Gr}
\DeclareMathOperator{\green}{\textit{a}}
\DeclareMathOperator{\HH}{H}
\DeclareMathOperator{\HHH}{\mathbb{H}}
\DeclareMathOperator{\Hom}{Hom}
\DeclareMathOperator{\id}{id}
\DeclareMathOperator{\Image}{Im}
\DeclareMathOperator{\kept}{\lambda}
\DeclareMathOperator{\Ker}{Ker}
\DeclareMathOperator{\len}{\ell}
\DeclareMathOperator{\py}{py}
\DeclareMathOperator{\Pic}{Pic}
\DeclareMathOperator{\Proj}{Proj}
\DeclareMathOperator{\pr}{p}
\DeclareMathOperator{\qp}{qp}
\DeclareMathOperator{\qr}{qr}
\DeclareMathOperator{\rank}{rank}
\DeclareMathOperator{\reg}{reg}
\DeclareMathOperator{\Sos}{\Sigma}
\DeclareMathOperator{\Span}{Span}
\DeclareMathOperator{\Spec}{Spec}
\DeclareMathOperator{\sw}{sw}
\DeclareMathOperator{\Sym}{Sym}
\DeclareMathOperator{\topint}{int}
\DeclareMathOperator{\Tor}{Tor}
\DeclareMathOperator{\tw}{tw}
\DeclareMathOperator{\variety}{V}
%% Personalized comments
\newcommand\mv[1]{{\color{blue}(MV: #1)}}
\newcommand\mve[1]{{\color{red}(MV: #1)}}
\newcommand\jh[1]{{\color{red}(JH: #1)}}

\title{Guía de estudio I (Computación Intermedia)}

\begin{document}
\maketitle
El objetivo de este documento es servir de guía para el estudiante en preparación para los parciales del curso.
Los enunciados marcados con Teorema son fundamentales y hay que conocer sus demostraciones.

\begin{enumerate}
\item \emph{Fundamentos de Álgebra Lineal:}
\begin{enumerate}
\item \emph{Teorema espectral (en dim finita)}

\begin{theorem}Toda matriz simétrica es ortogonalmente diagonalizable (sobre $\RR$).\end{theorem}
Más generalmente toda matriz normal ($A\in \CC^{n\times n}$ es normal $AA^*=A^*A$) es unitariamente diagonalizable (sobre $\CC$).
\item \emph{Descomposición en valores singulares:}
\begin{theorem} Si $A\in \RR^{m\times n}$ entonces existen matrices ortogonales $U\in \RR^{m\times m}$ y $V\in \RR^{n\times n}$ y una matriz $\Sigma$ diagonal con entradas reales tales que:
\[A=U\Sigma V^T\]
\end{theorem}
M\'as precisamente, podemos asumir que las entradas diagonales no nulas de $\Sigma$ son n\'umeros $\sigma_1\geq \dots \geq  \sigma_r>0$ donde $r={\rm rank}(A)$. Estos números se llaman los {\bf valores singulares} de $A$. 

Observaciones:

\begin{itemize}
\item La demostración de existencia es constructiva (si pudiéramos calcular valores y vectores propios, lo que haremos más adelante) pues se reduce a que:
\begin{enumerate}
\item Las columnas de $U$ pueden ser una base ortonormal de vectores propios de $AA^T$ (con $AA^Tu_i=\sigma_i^2u_i$), extendida a una base ortonormal del espacio completo si el rango lo requiere.
\item Las columnas de $V$ pueden ser una base ortonormal de vectores propios de $A^TA$ (con $A^tAv_i=\sigma_i^2v_i$), extendida a una base ortonormal del espacio completo si el rango lo requiere.
\item Los valores singulares son las raíces de los valores propios no nulos de $AA^T$ \'o $A^TA$. Esto demuestra la unicidad de los valores singulares y tambien explica en qué sentido falla la unicidad de $U$ y $V$.
\end{enumerate}
\item $\Sigma$ es, en general una matrix rectangular del mismo formato que $A$ y sus únicas entradas no nulas estan sobre la diagonal.
\item La descomposición de arriba permite, sumando sobre las entradas diagonales, escribir 
\[A=\sum_{i=1}^r \sigma_i u_iv_i^t\]
expresando $A$ como suma de $r$ matrices de rango uno.
\item Todo lo anterior es válido sobre $\CC$, donde $A\in \CC^{m\times n}$ y las matrices $U$ y $V$ son unitarias, llevando a una descomposición 
\[A=U\Sigma V^*\]
\end{itemize}
\end{enumerate}

\item \emph{Normas:}
\begin{enumerate}
\item Conocer el enunciado de la desigualdad de Holder y poder usarla para demostrar que $\|\bullet\|_p$ es norma en $\RR^n$ para $p\geq 1$. Saber demostrar:

\begin{theorem}
En un espacio vectorial de dimensión finita, todo par de normas son equivalentes.
\end{theorem}


\item \emph{Normas inducidas en espacios de matrices} Definición de la norma inducida sobre operadores $T:V\rightarrow W$ por normas $\|\cdot\|_V$ y $\|\cdot\|_W$ en $V$ y $W$
\[\|T\|_{\rm op}:=\sup_{v\in V\setminus\{0\}} \frac{\|Tv\|_W}{\|v\|_V}.\] 
Saber demostrar que es norma y que cumple $\|Tv\|_W\leq \|T\|_{\rm op}\|v\|_V$. 

\item Saber calcular $\|T\|_{\rm op,1,1}$, $\|T\|_{\rm op,\infty,\infty}$ en términos de las entradas de $T$ y expresar $\|T\|_{\rm op,2,2}$ en términos de valores singulares.  
\end{enumerate}


\item \emph{Números de condición y backward stability:}

\begin{enumerate}
\item Saber la definición formal de {\bf número de condición} y de {\bf número de condición absoluto} en un punto $x\in V$ de una función $f:V\rightarrow W$ entre espacios normados $V$ y $W$ (denotados $k(x,f)$ y $\tilde{k}(x,f)$ respectivamente).

\item Poder demostrar formalmente que el número de condición de {\it resolver el sistema lineal} $Ax=b$ para $b$ fijo esta acotado, para una matriz invertible $A$ por $k(A):=\|A\|\|A^{-1}|$. Si $\|\bullet\|:=\|\bullet\|_{\rm op,2,2}$ poder expresar $k(A)$ en términos de los valores singulares de $A$.

\item Entender la definición de backward-stable:
Para $\epsilon_0>0$ decimos que $\tilde{f}:V\rightarrow W$ es $\epsilon_0$-backward-stable con respecto a $f:V\rightarrow W$ si 
$\forall x\in V \exists \tilde{x}$ que cumple:
\begin{enumerate}
\item $\tilde{f}(x)=f(\tilde{x})$ 
\item $\|x-\tilde{x}\|\leq \epsilon_0\|x\|$
\end{enumerate}

\item Poder demostrar:
\begin{enumerate}
\item  Que una implementación del producto de números $x\ast y$ que garantice 
\[x\ast y = xy(1+\epsilon_0)\]
para todos $x,y$ es $\epsilon_0$-backward stable.
\item Que el producto interno de vectores en $\RR^n$, implementado con operaciones de producto y suma como las anteriores es backward-stable para algún múltiplo explícito de $\epsilon_0$ (para todo $\epsilon_0$ suficientemente pequeño). 

\item \emph{Uso de backward stability:}
\begin{theorem} Si $\tilde{f}$ es una implementación $\epsilon_0$-backward stable de $f$ entonces el error relativo en $x$ esta controlado por el número de condición $K(x,f)$  de $f$ en $x$. Más precisamente, si $f(x)\neq 0$ tenemos la siguiente desigualdad
\[\frac{\|\tilde{f}(x)-f(x)\|}{\|f(x)\|}\leq K(x,f)\epsilon_0\]
\end{theorem}

\item \begin{theorem} Resolver un sistema lineal triangular inferior mediante sustitución hacia atrás es backward-stable.
\end{theorem} 
Saber la demostración para sistemas $2\times 2$ (la idea en general es la misma y solo requiere más cálculos)

\end{enumerate}


\end{enumerate}
\item \emph{Factorización LU}:
\begin{enumerate}
\item Saber qué es la factorización $LU$ de una matriz cuadrada $A$.
\item Cómo usarías una factorización $A=LU$ conocida para resolver el sistema $Ax=b$? Qué puedes decir del número de condición y del error relativo de ese algoritmo?
\item Demostrar que si $A$ es invertible y la factorización $LU$ existe entonces es única. Existe siempre tal factorización? De existir es siempre única (removiendo el requisito de invertibilidad de $A$)?

\item Calcular la factorización $LU$ de la matriz 
$\begin{pmatrix}
\epsilon & 1 \\
1 & 1
\end{pmatrix}$
para $\epsilon\neq 0$. Notar la descomposición $LU$ no existe cuando $\epsilon=0$. Qué pasa con los números de condición de los factores $L$ y $U$ cuando $\epsilon\rightarrow 0$?


\item \emph{Criterio de existencia útil, sin pivoteo y consecuencias en matrices simétricas} Sea $A\in \RR^{n\times n}$ una matriz y sea $A_{\leq p}$ la submatriz cuadrada con índices de filas y columnas $\leq p$.

\begin{theorem} Si $\det(A_p)\neq 0$ para todo $p$ entonces existe una matriz triangular inferior con diagonal de $1$'s $L$ y una matriz triangular superior $U$ tal que $A=LU$.
\end{theorem}

\begin{theorem} [Sylvester + factorización de Cholesky] Si $A$ es simétrica y $\det(A_p)>0$ para todo $p$ entonces existe una única matriz invertible, triangular inferior $C$ con entradas positivas en la diagonal que cumple $A=CC^T$.
\end{theorem}

Más generalmente, si $\det(A_p)\neq 0$ para todo $p$ y $A$ es simétrica entonces existe una descomposición $A=LDL^T$ donde $L$ es triangular inferior con diagonal de unos y $D$ es una matriz diagonal.

\item Saber que, por Gauss-Jordan, TODA matriz admite una descomposición $LU$ con pivoteo $PA=LU$ para alguna permutación $P$. \emph{Dato curioso (sin dem.)} para TODA matriz simétrica $A$ existe una permutación $P$ tal que $P^T A P = LDL^T$ donde $L$ es triangular inferior con unos en la diagonal, $D$ es diagonal \emph{por bloques} donde todos los bloques tienen tamaño $2\times 2$ ó $1\times 1$). Ejemplo: $\begin{pmatrix}
0 & 1 \\
1 & 0
\end{pmatrix}$ necesita bloques $2\times 2$.



\end{enumerate}

\end{enumerate}

\end{document}
